% method_hmm.tex
% Subsection 3.1 body. Four narrative paragraphs: model and state
% partition; emissions and counted transitions; jump mechanism and
% initialization; and hyperparameter calibration. The bridge to the
% multi-asset composition lives at the top of method_sim.tex.


For each asset, the hidden Markov model with jumps (HMM-WJ) converted
the continuous return series into a
sequence of $N$ ordered return regimes. At each time step, the model
first selected a regime and then drew a return from that regime's
distribution (Fig.~\ref{fig:model_architecture}). We represented the model as
$\mathcal{M} = (\mathcal{S}, \mathcal{O}, \mathbf{T}, \mathbf{E}, \bar{\pi})$
~\cite{rabinerJuang1986}.
The hidden state $S_t \in \mathcal{S} = \{1,\dots,N\}$ identifies the
regime at time $t$, and the observed value is the excess growth rate
$G_t$ from Eq.~\eqref{eq:growth_rate}. We omit the ticker index
$i$ when discussing one asset. The transition matrix $\mathbf{T}$
sets the probabilities of moving between regimes, $\mathbf{E}$ gives
the return distribution within each regime, and $\bar{\pi}$ gives the
initial regime probabilities. We ordered the regimes from the most
negative to the most positive
returns. Their boundaries were equally spaced quantiles of a Laplace
distribution fitted to the asset's in-sample returns. Specifically,
$Q_k = F_L^{-1}(k/N;\, \mu_L, b_L)$, where $\mu_L$ and $b_L$ are the
fitted location and scale. We set $Q_0=-\infty$ and $Q_N=+\infty$ and
assigned return $G_t$ to state $k$ when $Q_{k-1}<G_t\le Q_k$. The
Laplace distribution captured the sharp concentration of small price
movements in Fig.~\ref{fig:empirical_motivation}a
~\cite{kotzKozubowskiPodgorski2001, tothJones2019}. We computed its
closed-form quantiles for all $424$ tickers~\cite{bilmes1998}.

The within-state emission followed a location-scale Student-$t$
distribution.
\begin{equation}
    G_t \mid S_t = k \;\sim\; \mu_k + \sigma_k \cdot t_{\nu},
    \qquad \nu = 5,
    \label{eq:emission}
\end{equation}
where $t_\nu$ denotes a standard Student-$t$ random variable with
$\nu$ degrees of freedom. For observations assigned to state $k$,
$\mu_k$ is the sample mean and $\sigma_k$ is the sample standard
deviation. A standard $t_\nu$ has variance $\nu/(\nu-2)$, so
$\sigma_k$ is a scale parameter rather than the emission standard
deviation. At $\nu=5$, the emission standard deviation is
$\sigma_k\sqrt{5/3}$. We selected $\nu = 5$ from a sensitivity sweep.
We used $\nu = \infty$ for the Normal baseline; Table~\ref{tab:hyperparams}
gives the explored range. Holding the rest of the model fixed, the
Student-$t$ emission raised simulated kurtosis toward the observed
value and raised the Kolmogorov--Smirnov (KS) and Anderson--Darling
(AD) pass rates relative to the Normal emission. We also
compared HMM-WJ with a hidden semi-Markov model (HSMM) baseline
adapted from Bulla and Bulla~\cite{bullaBulla2006}
(Table~\ref{tab:emission_hsmm}; Online
Appendix~\ref{sec:supp-emission-hsmm}). The emission and HSMM
comparisons separated the effects of heavy-tailed emissions and the
jump-duration mechanism. We estimated the transition matrix $\mathbf{T}$
by counting observed transitions, which avoided the initialization
sensitivity of iterative
expectation-maximization~\cite{bilmes1998}. The counted matrix was
concentrated near its diagonal, but individual states, including tail
states, were short-lived
(Fig.~\ref{fig:model_internals} in Online
Appendix~\ref{sec:supp-transitions}). The counted transition
probabilities caused the model to leave extreme-return states too
quickly to reproduce the observed
volatility clustering~\cite{bullaBulla2006}: a chain fitted by
counting reproduces the one-step state-pair frequencies exactly, and
its absolute-return autocorrelation therefore decays geometrically at
a rate set by the counted matrix. Online
Appendix~\ref{sec:supp-generator} derives the fitted generator's
stationary law and autocorrelation structure and quantifies both this
rapid mixing and the correction added by the jump mechanism.

We added jump episodes to extend visits to extreme-return states. When
no jump episode was active, a new episode began with probability
$\epsilon$. Its duration was
$K\sim\mathrm{Poisson}(\lambda)$, where $\lambda$ is the mean episode
length~\cite{glasserman2003}. Because a Poisson draw can be zero,
$K=0$ caused no forced tail step. For $K>0$, the model temporarily
replaced ordinary transitions with draws from two tail sets. The set
$\mathcal{S}_{-}=\{1,\dots,N_{\rm tail}\}$ contained the most negative
states, and $\mathcal{S}_{+}=\{N-N_{\rm tail}+1,\dots,N\}$ contained the
most positive states. At each forced step, the model selected the negative
set with probability $p_{\rm neg}$ and the positive set otherwise. It
made the tail-set choice independently at each step, so negative extreme
returns occurred slightly more often than positive ones.
Table~\ref{tab:hyperparams} lists $\epsilon$, $\lambda$, $p_{\rm neg}$,
and $N_{\rm tail}$. Two closed-form quantities summarize the
mechanism's operating point: forced steps occupy a long-run fraction
of approximately $\epsilon\lambda/(1+\epsilon\lambda)$ of all steps,
and a path of $T$ steps contains at least one episode with
probability approximately $1 - e^{-\epsilon T}$, so $\epsilon$ sets
how many paths jump while $\lambda$ sets how long an episode lasts
(Online Appendix~\ref{sec:supp-generator}). We obtained the
initial-state distribution
$\bar{\pi}$ by solving
$\bar{\pi}\mathbf{T}=\bar{\pi}$ with
$\sum_k\bar{\pi}_k=1$~\cite{hamilton2018regime, grinsteadSnell1997}.
Because $\mathbf{T}$ was estimated by counting, $\bar{\pi}$ equals
the vector of in-sample state frequencies up to end effects, so the
model's stationary growth-rate distribution is the
frequency-weighted mixture of the state emission distributions. The
solved $\bar{\pi}$ is exact for the ordinary transition matrix but only
approximate for the jump-augmented process because it does not include
the remaining jump duration. The full simulation procedure appears in
Algorithm~\ref{alg:hmmwj}. Each step followed $\mathbf{T}$ or continued
a jump episode and then drew a return from the selected state's emission
distribution.

% --- Architecture figure ----------------------------------------------------
% Main figure fig:model_architecture is collected after the references.

\begin{algorithm}[t]
\caption{Hidden Markov model with jumps (HMM-WJ) single-asset marginal generator.}
\label{alg:hmmwj}
\begin{algorithmic}[1]
  \REQUIRE Fitted parameters
    $(\mathbf{T}, \mathbf{E}, \bar{\pi}, \epsilon, \lambda,
    p_{\rm neg}, N_{\rm tail})$; horizon $T$;
    tail sets $\mathcal{S}_{-} = \{1, \dots, N_{\rm tail}\}$ and
    $\mathcal{S}_{+} = \{N - N_{\rm tail} + 1, \dots, N\}$.
  \ENSURE State sequence $\{S_t\}_{t=1}^{T}$ and growth-rate path
    $\{G_t\}_{t=1}^{T}$.
  \STATE Draw $S_1 \sim \bar{\pi}$;\quad
    draw $G_1 \sim \mu_{S_1} + \sigma_{S_1}\, t_{5}$;\quad
    $k \gets 0$
    \hfill \COMMENT{stationary start; $k$: remaining forced-step counter}
  \FOR{$t = 2, \dots, T$}
    \IF{$k > 0$}
      \STATE $\mathcal{J} \gets \mathcal{S}_{-}$ w.p.\ $p_{\rm neg}$, else $\mathcal{S}_{+}$;\quad
        $S_t \sim \mathrm{Uniform}(\mathcal{J})$;\quad $k \gets k - 1$
        \hfill \COMMENT{forced tail step; sign drawn per step}
    \ELSIF{$\mathrm{Uniform}(0,1) > \epsilon$}
      \STATE $S_t \sim \mathbf{T}_{S_{t-1}, \cdot}$
        \hfill \COMMENT{standard Markov transition}
    \ELSE
      \STATE $K \sim \mathrm{Poisson}(\lambda)$
        \hfill \COMMENT{trigger jump; $K$ is the episode length}
      \IF{$K = 0$}
        \STATE $S_t \sim \mathbf{T}_{S_{t-1}, \cdot}$
          \hfill \COMMENT{$K = 0$: no forced steps, standard transition}
      \ELSE
        \STATE $\mathcal{J} \gets \mathcal{S}_{-}$ w.p.\ $p_{\rm neg}$, else $\mathcal{S}_{+}$;\quad
          $S_t \sim \mathrm{Uniform}(\mathcal{J})$;\quad
          $k \gets \min(K - 1, T - t)$
          \hfill \COMMENT{first forced step; $K$ steps total, sign per step}
      \ENDIF
    \ENDIF
    \STATE Draw $G_t \sim \mu_{S_t} + \sigma_{S_t}\, t_{5}$
      \hfill \COMMENT{Eq.~\eqref{eq:emission}}
  \ENDFOR
\end{algorithmic}
\end{algorithm}

We chose the jump probability $\epsilon$ and mean duration $\lambda$
to match the SPY absolute-return autocorrelation and
kurtosis~\cite{mandelbrot1963, cont2001,
rydenTerasvirtaAsbrink1998}. Jump-diffusion models use Poisson events to
add instantaneous price jumps~\cite{merton1976, kou2002}. In HMM-WJ, a
jump episode extended a visit to the tail states. We searched a grid of
$(\epsilon,\lambda)$ values. At each grid point, we averaged squared
errors in the autocorrelation function (ACF) of absolute returns through lag $L$ and
weighted squared errors in excess kurtosis over the jump-active paths.
We minimized the objective $J(\epsilon,\lambda)$.
\begin{equation}
    J(\epsilon, \lambda) =
    \frac{1}{|\mathcal{R}_J|}\sum_{r\in\mathcal{R}_J}
    \left[
    \sum_{\tau=1}^{L} \left( \mathrm{ACF}_{\rm obs}(\tau)
    - \mathrm{ACF}_{r}(\tau) \right)^2
    + w_\kappa \left( \kappa_{\rm obs} - \kappa_{r} \right)^2
    \right],
    \label{eq:grid_search}
\end{equation}
where $\mathcal{R}_J$ is the set of jump-active simulated paths,
$\mathrm{ACF}_{\rm obs}(\tau)$ and $\kappa_{\rm obs}$ are the empirical
absolute-growth autocorrelation at lag $\tau$ and excess kurtosis, and
$\mathrm{ACF}_{r}(\tau)$ and $\kappa_r$ are the corresponding values
for path $r$. We used $\kappa$ for kurtosis and $K$ for jump duration.
We simulated $200$ independent paths of $2{,}766$ trading days at
each grid point. A path with no jump reduces to the hidden Markov model
with no jumps (HMM-NJ) and contains no
information about $(\epsilon,\lambda)$, so the objective included only
paths with at least one jump~\cite{glasserman2003}. Within those
paths, $\lambda$ controls both the height and the horizon of the slow
autocorrelation component while the kurtosis penalty limits the total
forced-step mass; $\epsilon$ acts mainly through the share of paths
that contain an episode
(Online Appendix~\ref{sec:supp-generator}). The value of
$w_\kappa$ set the relative weight of the kurtosis error. The search ranges,
$w_\kappa$, $L$, and the remaining HMM-WJ and SIM hyperparameters are
reported in Table~\ref{tab:hyperparams}. We fixed the selected
hyperparameters before running the generator comparisons and the
holdout evaluation.
